% ============================================================
\section{Giới thiệu}
\label{sec:intro}
% ============================================================

\subsection{Bối cảnh và mục tiêu}

Bài tập lớn yêu cầu xây dựng một \textbf{trợ lý hỗ trợ học tập} (Course Assistant) cho
môn Xử lý Ngôn ngữ Tự nhiên (CO3085): người dùng đặt câu hỏi tiếng Việt về thông tin
môn học, lịch học theo tuần, chương và chủ đề, chuẩn đầu ra (LO), bài tập lớn,
deadline, tài liệu tham khảo và quy định của môn; hệ thống phải trả lời dựa trên
knowledge base (KB) do giảng viên cung cấp \cite{spec}.

Điểm cốt lõi của đề bài không nằm ở việc ``trả lời được'' mà ở việc thể hiện đầy đủ
một \textbf{pipeline NLP} có thể giải thích từng bước:
\begin{center}
  \small
  Câu hỏi $\rightarrow$ Grammar/Parser $\rightarrow$ Cây cú pháp $\rightarrow$
  Diễn giải ngữ nghĩa $\rightarrow$ Intent/Entity $\rightarrow$ Truy vấn KB
  $\rightarrow$ Câu trả lời
\end{center}
Khi KB không có dữ kiện, hệ thống phải báo rõ là không tìm thấy thay vì tự tạo câu
trả lời.

\subsection{Phạm vi và hướng tiếp cận}

Nhóm chọn \textbf{Phần II -- Lựa chọn 1: Classical NLP}. Toàn bộ xử lý dựa trên văn
phạm phi ngữ cảnh (CFG) viết tay, bộ phân tích cú pháp Earley và luật diễn giải cây
cú pháp thành biểu diễn ngữ nghĩa. Chương trình chạy offline bằng thư viện chuẩn
Python, không gọi API hay mô hình ngôn ngữ. Lý do lựa chọn:
\begin{itemize}
  \item Miền hỏi đáp hẹp và đóng (một môn học, sáu file KB), phù hợp với văn phạm
        ngữ nghĩa (semantic grammar) như nội dung các chương 3--5 của môn học.
  \item Mọi câu trả lời lần ngược được về truy vấn, frame ngữ nghĩa và cây cú pháp;
        kết quả tất định nên dễ kiểm thử và chấm theo từng tầng.
  \item Phân biệt rõ hai loại thất bại: câu nằm ngoài văn phạm (không có cây) và câu
        hiểu được nhưng KB thiếu dữ kiện (có cây, có intent, truy vấn rỗng).
\end{itemize}

\Cref{tab:scope} tóm tắt mức độ thực hiện từng phần của đề và file output tương ứng.

\begin{table}[htbp]
  \centering
  \small
  \begin{tabularx}{\linewidth}{|P{2.6cm}|P{1.9cm}|L|P{3.3cm}|}
    \hline
    \textbf{Phần} & \textbf{Yêu cầu} & \textbf{Hiện thực} & \textbf{Output} \\ \hline
    I. Văn phạm và parser & Bắt buộc & CFG 26 nhánh câu hỏi, parser Earley, bộ sinh
      câu có kiểm tra parse lại & \code{grammar.txt}, \code{samples.txt},
      \code{parse-results.txt} \\ \hline
    II. Ngữ nghĩa và hỏi đáp & Bắt buộc & Classical: luật diễn giải cây, entity lexicon,
      truy vấn KB, sinh câu trả lời theo template & \code{semantic.txt},
      \code{intent-entity.txt}, \code{query.txt}, \code{answer.txt} \\ \hline
    III. Ngữ cảnh cục bộ & Mở rộng & Bộ nhớ thực thể gần nhất theo loại, 5 dạng tham
      chiếu & \code{dialogue.txt} \\ \hline
    IV. Đánh giá & Bắt buộc & Chấm end-to-end trên 42 câu của đề, 22 câu thử thách và
      6 hội thoại & \code{evaluation.txt} \\ \hline
  \end{tabularx}
  \caption{Phạm vi thực hiện theo các phần của đề}
  \label{tab:scope}
\end{table}

\subsection{Thành viên và phân công}

Bảng~\ref{tab:members} tổng hợp phân công theo lịch sử commit và pull request của
nhóm. Giai đoạn 1, mỗi thành viên viết grammar cho một nhóm nhánh câu hỏi; giai
đoạn 2, mỗi người hiện thực các tầng của pipeline theo hợp đồng chung
(\cref{sec:arch}).

\begin{table}[htbp]
  \centering
  \small
  \begin{tabularx}{\linewidth}{|P{3.6cm}|P{1.7cm}|L|}
    \hline
    \textbf{Thành viên} & \textbf{MSSV} & \textbf{Công việc chính} \\ \hline
    Phan Gia Bảo & 2310251 & Grammar 9 nhánh thông tin môn học, giảng viên, lịch theo
      tuần/chương; tokenizer, parser Earley, bộ sinh câu. \\ \hline
    Nguyễn Huy Bách & 2310198 & Grammar 8 nhánh chủ đề, LO, bài tập lớn; nạp KB
      (\code{kb.py}) và bộ truy vấn (\code{query.py}). \\ \hline
    Mai Anh Khoa & 2311601 & Grammar 9 nhánh deadline, tài liệu, quy định và câu ngữ
      cảnh; diễn giải ngữ nghĩa, sinh câu trả lời, ngữ cảnh hội thoại. \\ \hline
    Hoàng Ngô Thiên Phúc & 2212612 & Review và gộp grammar; hợp đồng giữa các tầng,
      golden test; pipeline, CLI, xuất output, đánh giá; mở rộng grammar; Docker,
      README, báo cáo, video demo, đóng gói. \\ \hline
  \end{tabularx}
  \caption{Thành viên và phân công}
  \label{tab:members}
\end{table}

\subsection{Cấu trúc báo cáo}

\Cref{sec:data} mô tả dữ liệu và KB. \Cref{sec:arch} trình bày kiến trúc và ranh giới
giữa các tầng. \Cref{sec:part1,sec:part2,sec:part3} lần lượt đi qua Phần I (grammar,
parser, sinh câu), Phần II (ngữ nghĩa, truy vấn, trả lời) và Phần III (ngữ cảnh).
\Cref{sec:eval} báo cáo kết quả đánh giá và phân tích lỗi; \cref{sec:discussion} nêu
ưu điểm, hạn chế và hướng mở rộng. \Cref{sec:reproduce} hướng dẫn chạy lại; cuối cùng
là nguồn tham khảo và khai báo sử dụng AI.
